\section{The finite extension searches}\label{app:search}

\subsection{Enumeration of the extension maps}

For each classified shortened code, the following procedure is applied.
The local quotient has a fixed physical symplectic basis of size \(2\mu\).
The search first chooses the images of one isotropic half of this basis,
then chooses the images of its symplectic partners. In matrix notation
these two groups of columns are \(X\) and \(Y\). A partial assignment
is discarded only when it already violates a necessary condition; every
assignment that survives is extended through all remaining allowed values.
\begin{enumerate}[label=\arabic*.]
\item Enumerate the permitted local kernels \(K\), shortened-code
kernels \(N\), and radical identifications, retaining representatives
under the physical groups specified below.
\item Assign the \(\mu\) images \(X\) one at a time. Test all newly
determined linear combinations, the kernel evaluations, and the
necessary partial-rank and weight-count conditions.
\item For a complete \(X\), transform the Gram matrix as in
\eqref{eq:Y-normal}. Reject it if the lower-right block is nonzero.
\item Enumerate the remaining free entries of \(U\), including its
diagonal. Test the kernel equations and every complementary mask
against its full distance domain \(\mathcal D_x\).
\item Reconstruct each completed map in the original physical coordinates
and test the required stabilizer weight counts. Save any survivor.
\end{enumerate}
The sixteen-qubit relaxation omits the Gram and prescribed-kernel tests
and assigns all basis images using only the distance inequalities.
Every genuine extension occurs in the classified domain, passes the
necessary partial tests, and is retained by the physical symmetry
reductions. If no completed assignment passes the final tests, no
genuine extension exists. A survivor of the sixteen-qubit relaxation
would require a separate commutation check.

\subsection{Coset distances from physical words}

Using a parity-check basis of \(C\), compute the syndrome of each word
\(w\in V_f\) and minimize \(\wt(w)\) for each syndrome.
Equivalently, construct the Cayley graph of \(W/C\) whose edges add
one nonzero Pauli at one site, and use breadth-first search from zero.
A path of length \(u\) gives a physical word of weight at most \(u\),
while a word of weight \(v\) gives a path of length \(v\). Therefore
the two procedures compute the same exact coset minimum. Rank checks
verify that the syndrome kernel is \(C\).

\subsection{Completing a partial extension}

Write \(\dim E=2\mu\), choose a symplectic basis, and write \(M=[X\ Y]\),
where \(X\) contains images of classes represented using only the letters
\(I\) and \(Z\). We call these the \(Z\)-only classes. Every nonzero \(Z\)-only
class has minimum weight below the target distance in the searches that use this
normal form. Its image cannot vanish, so \(X\) has rank \(\mu\).
Complete its columns to an invertible syndrome basis \(B_X\), and put
\[
 X'=B_X^{-1}X=\begin{bmatrix}I_\mu\\0\end{bmatrix},\quad
 Y'=B_X^{-1}Y=\begin{bmatrix}U\\L\end{bmatrix},\quad
 \Gamma'=B_X^{-1}\Gamma B_X^{-\tr}.
\]
The compatibility equation is
\begin{equation}\label{eq:Y-normal}
 \Gamma'=\begin{bmatrix}U+U^\tr&L^\tr\\L&0\end{bmatrix}.
\end{equation}
The lower-right block must vanish; \(L\) is fixed; and
\(U+U^\tr\) fixes each off-diagonal pair sum while leaving every
diagonal free. For \(\mu=6\), there are 21 binary free entries before
distance and kernel constraints.

For each nonzero complementary mask \(x\), form the domain
\[
 \mathcal D_x=\{u\in\F^\mu:\delta(B_X(z+u,Lx))+
                       c_R(z,x)\geq d\text{ for every }z\in\F^\mu\}.
\]
Here \(x\) is a binary coefficient vector selecting a sum of the
complementary basis vectors, while \(z\) selects a sum of the first
basis vectors. The matrix \(L\) is already fixed. For a proposed value
\(Ux=u\), the displayed test checks all words with those complementary
coefficients, by letting \(z\) run over \(\F^\mu\).
The remaining exact constraints are \(Ux\in\mathcal D_x\) for all
\(x\ne0\), together with \eqref{eq:Y-normal} and any prescribed
kernel equations. A finite traversal of all allowed entries of \(U\)
therefore enumerates every completion. Consistency tests on pairs of assignments prune the domains. The final
check uses every mask, including those of size greater than two.

For a fixed kernel vector with coordinates \((z,x)\), the additional
equation is \(Xz+Yx=0\); in this normal form it reads
\(Ux=z\), \(Lx=0\). The implementation retains the resulting affine
freedom and exhausts all completions for each retained assignment to
\(X\).

\subsection{Partial rank and exact profile tests}

After assigning \(j\) \(Z\)-only images, let \(\mathcal U_j\) be their common
annihilator in the coefficient space of \(P\). Restricting
\(\Gamma=XY^\tr+YX^\tr\) to \(\mathcal U_j\) makes the assigned terms vanish.
Only \(\mu-j\) pairs remain, so the restricted rank is at most
\(2(\mu-j)\): each remaining pair contributes a matrix
\(uv^\tr+vu^\tr\) of rank at most two. This is a necessary partial-depth test, with a zero
annihilator block at full depth.

We also use the weight distribution of the reconstructed stabilizer.
For a coefficient vector \(a\), write
\(v(a)=\sum a_iv_i\) and
\(e(a)=JM^\tr a=(Y^\tr a,X^\tr a)\). Its words are
\begin{equation}\label{eq:physical-S}
 (\widetilde e(a)+r,v(a)),\qquad a\in\F^b,\ r\in R.
\end{equation}
All pairs \((a,r)\) give distinct physical words, since the suffix basis
is independent. The words with \(a=0\) are \(R\).
This remains true when \(\Gamma\) is singular.

At full \(X\), for \(a\in\Ann(X)\), compatibility gives
\(\Gamma a=X(Y^\tr a)\). Independence of \(X\) determines
\(Y^\tr a\) uniquely when it exists. Every translate with such a coefficient can then be counted exactly.
The estimates below are used only for the other coefficients, so each
word is counted once. At completed \(M\),
all words in \eqref{eq:physical-S} are counted exactly and the required
profile equality is checked. No inversion of a singular \(\Gamma\)
is used.

For the remaining coefficients, we obtain lower bounds on the number
of weight-six and weight-seven words as follows.
For a fixed nonzero coefficient \(a\), let \(\mathcal P_a\) contain
every prefix representative that any completion of the current partial
map could assign to \(v(a)\). A nonempty superset is sufficient; if
there are no possible prefixes, the partial assignment has no completion.
Define
\[
 U_a=\wt(v(a))+\max_{v\in\mathcal P_a}\min_{r\in R}\wt(v+r).
\]
The order of the maximum and minimum is important. For each possible
prefix, first take its least weight modulo \(R\); then take the
largest value any completion could achieve. Even this best possible
completion cannot make the corresponding coset heavier than \(U_a\).
For every completion, the minimum weight in the corresponding
\(R\)-coset is at most \(U_a\). In the pure distance-six setting,
\(U_a\leq7\)
therefore guarantees at least one word of weight six or seven, and
\(U_a\leq6\) guarantees a weight-six word. Distinct coefficients
give disjoint cosets by \eqref{eq:physical-S}, so these guaranteed
counts may be added. If they exceed the allowed counts outside \(R\),
no completion exists. The same argument applies to bounds on individual
translates. Because the maximization uses a superset of the possible
prefixes, it can weaken the rejection test but cannot remove a valid
completion. Exact profile equalities are tested only after the whole
map is assigned.

\subsection{Exact prescribed kernels}\label{app:fixedkernel}

For fixed \(N\subseteq\rad(P)\), work in \(\bar P=P/N\).
Let its alternating matrix be \(\bar\Gamma\). A genuine extension
gives an injective isometry \(\tau:\bar P\to K^\perp\).
Let \(g=q-t\); fix every isomorphism from \(\rad(\bar P)\)
to \(\rad(K)\). Relative to chosen bases, this specifies evaluations
of the \(Z\)-only image columns against every radical coefficient.
The coordinate labels given by these \(g\) radical evaluations span
\(\F^g\). A site permutation makes the first \(g\) labels
independent. The same permutation is applied to the local quotient
basis and its table of coset minimum weights.

Subtract suitable combinations of those first \(g\) image columns
to obtain vectors \(T_i\in\im\bar\Gamma\), and make the same
subtractions from the physical prefix vectors. When \(t=2\), project
the latter off a fixed symplectic pair spanning the nondegenerate
part of \(K\). Call the resulting vectors \(e_i\).
For any solution \(c_i\) of \(\bar\Gamma c_i=T_i\), compatibility
requires
\[
 c_i^\tr T_j=\langle e_i,e_j\rangle.
\]
Changing \(c_i\) by a radical vector does not alter the left-hand
side. These equations, all radical evaluations, and all \(Z\)-only
weight inequalities are necessary conditions during the search over
\(X\). Every retained assignment is then tested by the complete search
over \(Y\).

The final completion enforces \(M(K)=0\) and the full Gram equation.
The prescribed radical identification also fixes the values of
\(X^\tr\) on \(\rad(P)\). Its kernel there is \(N\):
on \(\rad(P)/N\) it is the identification with \(\rad(K)\),
followed by projection onto the second half of the local coordinates.
That projection is injective on \(K\), because a nonzero element in
its kernel would be a \(Z\)-only class of weight less than \(d\).
The following argument shows why these conditions give the prescribed
kernels exactly.

Identify \(P\) with its coefficient space \(\F^b\), so its radical
is \(\ker\Gamma\), and write \(\tau=JM^\tr\).
If \(\tau c=0\), the Gram equation gives \(\Gamma c=0\);
also \(X^\tr c=0\). The prescribed radical evaluations therefore
give \(c\in N\). Conversely, if \(c\in N\), then
\(\Gamma c=0\) and \(X^\tr c=0\). Expanding the Gram equation,
\[
 0=\Gamma c=X(Y^\tr c)+Y(X^\tr c)=X(Y^\tr c).
\]
The columns of \(X\) are independent, so \(Y^\tr c=0\).
\(\tau c=0\), proving \(\ker\tau=N\). Finally,
\[
 \rank M=\rank\tau=b-\dim N=2\mu-q,
 \qquad \dim\ker M=q.
\]
Since \(K\subseteq\ker M\) and \(\dim K=q\), we obtain
\(\ker M=K\). Requiring containment of the prescribed subspaces
therefore suffices here because the radical conditions also fix the
dimensions; larger kernels cannot pass these tests.

\subsection{Physical symmetry and joint normalization}\label{app:symmetry}

The symmetry reductions use site permutations and local permutations
of the nonzero Pauli letters. These preserve both the symplectic form
and the minimum weights in the original coordinates.

A physical suffix automorphism acts on coefficient vectors by \(G_{\rm coef}\):
\(g(v(a))=v(G_{\rm coef}a)\). Its syndrome action is \(G_{\rm syn}=G_{\rm coef}^{-\tr}\), since
\(\sigma(gw)=G_{\rm coef}^{-\tr}\sigma(w)\). \(M\mapsto G_{\rm syn}M\), and
\(G_{\rm syn}\Gamma G_{\rm syn}^\tr=\Gamma\). This contragredient action is essential
when \(\Gamma\) is singular. Every generator is checked on physical
words, coset distances, and the corresponding form.

Outer prefix automorphisms may reduce radical identifications if they
preserve the local space and distinguished local stabilizer. Inside a
fixed identification, only actions fixing its radical classes are used.
Suffix actions likewise preserve the specified \(C,N\) and radical
identification. These restrictions ensure that both actions preserve
the set of extensions with the prescribed kernels and identification.

Choose a lexicographically least \emph{whole} \(Z\)-only tuple in a
joint orbit of those prefix and suffix actions. It is lexicographically minimal under the prefix action,
and its first label is least in its suffix orbit. Otherwise one of the
actions would reduce the same whole tuple. Imposing both conditions
therefore retains a representative of every joint orbit, even when the
two actions do not commute. No uniqueness of a joint canonical form
is required.

For the mixed six-site permutation reduction, the fixed first image is
unchanged and the other images are differences from an anchor among six
points. Physical permutations reorder and reanchor these points. A
lexicographically least full tuple is increasing and no translated
anchor tuple is smaller. The search assigns points in increasing order,
so its current tuple is a fixed initial segment of every completion of
that branch. Under a change of anchor, adding future points can only
decrease the initial entries of the sorted translated list. If the
already assigned points give a smaller initial segment after translation,
every completion has a physically equivalent tuple smaller
than itself. Such a branch cannot contain the least representative and
can be discarded.
The same physical action is transported to the complementary basis,
including its possible shear. Optional relaxed complementary-domain
filters that are not invariant under reanchoring are disabled in that
search; the complete \(Y\)-test is retained.

The full-plane eight-site reduction uses the corresponding physical
eight-point action and verifies it on the full quotient basis.
All first-label orbit lists explicitly record every eligible label and
its orbit membership. Each list is divided into consecutive blocks
of at most 64 labels. Searching every block covers every
eligible first image.
